\section{Инструменты исследования}

В рамках исследования качества программы применялся следующий стек утилит:
\begin{itemize}
    \item \textbf{gcov (GNU Coverage)} --- для проверки полноты покрытия кода тестами.
    \item \textbf{Valgrind Memcheck} --- для отслеживания утечек памяти и ошибок доступа.
    \item \textbf{Valgrind Massif} --- для построения профиля потребления кучи (heap memory) во времени и выявления пиковой нагрузки.
    \item \textbf{gprof (GNU Profiler)} --- для анализа процессорного времени (CPU time) и выявления «узких мест» (bottlenecks) в архитектуре программы.
\end{itemize}

Для тестирования был написан скрипт на Python, генерирующий файлы с 1,000,000 команд вставок и поисков (случайные данные и отсортированные данные).

\section{Исследование исходной (рекурсивной) версии}

Изначальная реализация АВЛ-дерева использовала рекурсивные алгоритмы спуска и вращений, а управление памятью осуществлялось посредством умных указателей \texttt{std::unique\_ptr}. 

\subsection*{Анализ памяти (Valgrind)}
Инструмент \texttt{memcheck} подтвердил абсолютную безопасность работы с памятью:
\begin{alltt}
==4521== HEAP SUMMARY:
==4521==     in use at exit: 0 bytes in 0 blocks
==4521== All heap blocks were freed -- no leaks are possible
\end{alltt}
Анализ пикового потребления памяти с помощью \texttt{Massif} показал пик в \textbf{68.83 MB} для 1 миллиона узлов. С учётом полезной нагрузки структуры узла (ключ, значение, баланс, два указателя), накладные расходы аллокатора составили около 8 байт на каждый вызов \texttt{new}, что является нормальным поведением для C++.

\subsection*{Анализ времени выполнения (gprof)}
Наибольший интерес представили результаты профилирования процессорного времени на отсортированном наборе данных (где АВЛ-дерево совершает максимум балансировок).

Выдержка из отчёта \texttt{gprof} до оптимизации:
\begin{alltt}
  %   cumulative   self              self     total           
 time   seconds   seconds    calls  ns/call  ns/call  name    
 11.71      0.33     0.13 81805701     1.59     8.11  frame_dummy
 10.81      0.45     0.12 363625137    0.33     1.05  std::__uniq_ptr_impl::_M_ptr()
  7.21      0.69     0.08 161611220    0.50     1.71  std::unique_ptr::operator bool()
  5.41      0.82     0.06 363625137    0.17     0.17  std::unique_ptr::get()
  3.60      0.91     0.04 61854256     0.65     1.25  std::unique_ptr::~unique_ptr()
  ...
  0.00      1.11     0.00  1000000     0.00   910.00  AVLTree::Insert()
\end{alltt}

В таблице отчётливо видно, что почти \textbf{80\% времени работы программы} занимает не сам алгоритм балансировки, а вызовы методов стандартной библиотеки для обслуживания \texttt{std::unique\_ptr} (\texttt{get}, \texttt{operator bool}, разыменование \texttt{tuple} внутри умного указателя). Рекурсия накладывает огромный штраф на раскрутку стека вызовов (\texttt{frame\_dummy}).

\textbf{Найденный недочёт:} Высокие накладные расходы C++ инфраструктуры (умные указатели и рекурсия) перекрывают алгоритмическую эффективность дерева поиска.
\pagebreak